\chapter{sotFS: el grafo Merkle como filesystem}
\label{cap:sotfs-modelo}

\section{Motivación}

Filesystems tradicionales (\emph{ext4}, \emph{XFS}, \emph{NTFS})
modelan los datos como un \emph{árbol}: directorios contienen
inodes, inodes apuntan a bloques. La estructura es jerárquica;
el mismo bloque pertenece a un solo archivo. Esa decisión tiene
costos:

\begin{itemize}
  \item \emph{Hard links} son un bolt-on: dos paths que apuntan
    al mismo inode existen, pero el modelo subyacente no los
    refleja como \emph{misma cosa con dos nombres}.
  \item Deduplicación post-escritura es cara: hay que rescanear
    bloques iguales y reemplazar referencias.
  \item Snapshots se hacen vía COW (\emph{copy on write}) o
    \emph{shadow paging}; son técnicas, no propiedades del
    modelo.
\end{itemize}

\texttt{sotFS} usa una estructura distinta: un \emph{Merkle DAG}
(grafo dirigido acíclico con direccionamiento por hash), donde
cada nodo se nombra por el hash criptográfico de su contenido
\cite{Merkle1987-DAG}. La estructura tiene:

\begin{itemize}
  \item \textbf{Deduplicación intrínseca.} Dos nodos con el mismo
    contenido tienen el mismo hash; son el mismo nodo.
  \item \textbf{Snapshots gratis.} Un snapshot es un puntero a
    la raíz del DAG en un momento; nada se copia.
  \item \textbf{Verificación cripto.} Dado un hash de raíz, el
    contenido completo del subgrafo se puede verificar
    incrementalmente.
  \item \textbf{Sharing entre dominios.} El mismo bloque puede
    referenciarse desde dos archivos sin bolt-on.
\end{itemize}

El precio: la estructura es un \emph{DAG con seis tipos de
nodo y seis tipos de arista}; las operaciones POSIX se modelan
como \emph{double pushout} (DPO) graph rewriting. La
verificación formal vive en \texttt{formal/sotfs\_graph.tla}.

\section{Los seis tipos de nodo}

\begin{lstlisting}[frame=single, basicstyle=\ttfamily\scriptsize]
Inode      — metadata (mode, uid, mtime); apunta a Blocks
Dir        — directorio; contiene entradas (Name → Inode)
Block      — bloque de datos (4 KiB); identificado por hash
Snapshot   — root pointer + timestamp
Capability — capability sotFS (subset de la cap del kernel)
Domain     — partición lógica del filesystem
\end{lstlisting}

Y los seis tipos de arista:

\begin{lstlisting}[frame=single, basicstyle=\ttfamily\scriptsize]
contains    Dir → Inode      (el directorio incluye este inode)
extent      Inode → Block    (este bloque pertenece al inode)
named-as    Dir → Name → Inode  (la entrada con este nombre)
parent      Inode → Dir      (inverso de contains; navegabilidad)
captured    Capability → Inode  (la cap autoriza este inode)
owns        Domain → Inode   (cuál dominio es dueño)
\end{lstlisting}

\section{Por qué DAG y no árbol}

\begin{figure}[h]
  \centering
  \begin{tikzpicture}[
      every node/.style={font=\scriptsize},
      ino/.style={node base, fill=sotxBlue!15, draw=sotxBlue!70,
                  minimum width=20mm},
      blk/.style={node base, fill=sotxGreen!15, draw=sotxGreen!70,
                  minimum width=18mm},
      arr/.style={-Stealth, draw=sotxFg, line width=0.5pt}
    ]
    % Inodes
    \node[ino] (i1) at (-3,2) {Inode A\\\texttt{file1.txt}};
    \node[ino] (i2) at (3,2)  {Inode B\\\texttt{file2.txt}};
    % Blocks
    \node[blk] (b1) at (-4,0) {Block\\hash=A1};
    \node[blk] (b2) at (0,0)  {Block\\hash=B2 (shared)};
    \node[blk] (b3) at (4,0)  {Block\\hash=C3};

    % Edges
    \draw[arr] (i1) -- (b1) node[font=\scriptsize\ttfamily,midway,fill=sotxBg]
       {extent};
    \draw[arr] (i1) -- (b2);
    \draw[arr] (i2) -- (b2);   % SHARED
    \draw[arr] (i2) -- (b3);

    \node[hazard, anchor=north west, font=\scriptsize\itshape]
      at (-4,-1.5)
      {Block B2 es compartido por dos inodes distintos sin\\
       hard-link explícito; emerge del Merkle deduplication.};
  \end{tikzpicture}
  \caption{Sharing intrínseco: \texttt{file1.txt} y
  \texttt{file2.txt} comparten un mismo bloque por construcción
  porque tienen el mismo hash. No es un \emph{hard link}; es la
  semántica natural del DAG.}
  \label{fig:sotfs-sharing}
\end{figure}

\section{Direccionamiento por hash}

Cada nodo tiene un \emph{content hash}: BLAKE3 de 32 bytes sobre
su contenido serializado. La función:

\begin{lstlisting}[language=rust, frame=single, basicstyle=\ttfamily\scriptsize]
fn node_hash(node: &Node) -> [u8; 32] {
    let mut hasher = Blake3::new();
    hasher.update(&node.type_tag.to_le_bytes());
    hasher.update(&node.serialize());
    hasher.finalize().into()
}
\end{lstlisting}

\begin{invariante}[Identidad por contenido]
\label{inv:sotfs-identity}
Dos nodos con el mismo contenido tienen el mismo hash; el
filesystem mantiene una sola copia. \emph{Identidad estructural
implica identidad de almacenamiento}.
\end{invariante}

\begin{invariante}[Acíclico]
\label{inv:sotfs-acyclic}
Los punteros entre nodos son siempre por hash, y un nodo no
puede contener su propio hash (los hashes son criptográficamente
fuertes; la probabilidad de un \emph{ciclo accidental} es
$2^{-256}$). El grafo es DAG por construcción.
\end{invariante}

\begin{invariante}[Inmutabilidad de nodo]
\label{inv:sotfs-immutable}
Un nodo nunca se modifica \emph{en sitio}. Una mutación lógica
(truncate, append, chmod) crea un nuevo nodo con el contenido
nuevo, deja el viejo si nadie lo referencia, e \emph{actualiza
el padre} apuntando al nuevo hijo. Las roots de snapshot
preservan la vista vieja.
\end{invariante}

\section{DPO: double pushout para operaciones POSIX}

Toda operación POSIX se modela como una transformación DPO:

\begin{enumerate}
  \item \textbf{Match.} El operador especifica el \emph{LHS}
    (subgrafo que tiene que existir antes).
  \item \textbf{Glue.} Un subgrafo que persiste tras la
    operación.
  \item \textbf{Result.} El \emph{RHS}, que reemplaza el LHS
    manteniendo el glue.
\end{enumerate}

Para \texttt{rename(d/old, d/new)}:

\begin{itemize}
  \item LHS: $d \xrightarrow{\text{named-as old}} i$.
  \item Glue: $d$ y $i$ permanecen.
  \item RHS: $d \xrightarrow{\text{named-as new}} i$.
\end{itemize}

Para \texttt{unlink(d/name)}:

\begin{itemize}
  \item LHS: $d \xrightarrow{\text{named-as name}} i$.
  \item Glue: $d$ permanece.
  \item RHS: $i$ se borra (si \texttt{i.refcount == 0}); la arista desaparece.
\end{itemize}

\section{El \texttt{.tla} de \texttt{sotfs\_graph}}

\begin{lstlisting}[language=tla, frame=single, basicstyle=\ttfamily\scriptsize]
VARIABLES
    inodes,         \* Set of active inode IDs
    dirs,           \* Set of active directory IDs
    blocks,         \* Set of active block IDs
    inodeType,      \* Function: InodeId -> {"regular", "directory"}
    inodeBlocks,    \* Function: InodeId -> SeqSet of BlockIds
    dirEntries      \* Function: DirId -> Set of <<name, InodeId>>

\* DPO rule: link an existing inode under a new name in a directory.
LinkInode(d, i, name) ==
    /\ d \in dirs
    /\ i \in inodes
    /\ ~\E e \in dirEntries[d] : e[1] = name      \* gluing condition
    /\ dirEntries' = [dirEntries EXCEPT
                       ![d] = @ \cup {<<name, i>>}]
    /\ UNCHANGED <<inodes, dirs, blocks, inodeType, inodeBlocks>>
\end{lstlisting}

La cláusula \texttt{\textasciitilde\textbackslash E e \textbackslash in
dirEntries[d] : e[1] = name} es la \emph{gluing condition}: la
operación está habilitada sólo si no existe ya una entrada con
ese nombre. Es la versión TLA del \texttt{EEXIST} de POSIX, pero
no como un \emph{check ad hoc}; emerge de la estructura DPO.

\section{Las invariantes verificadas}

\begin{lstlisting}[language=tla, frame=single, basicstyle=\ttfamily\scriptsize]
\* Toda entrada de directorio apunta a un inode vivo.
DirIntegrity ==
    \A d \in dirs : \A e \in dirEntries[d] : e[2] \in inodes

\* Todo bloque referenciado por un inode existe.
ExtentIntegrity ==
    \A i \in inodes : \A b \in inodeBlocks[i] : b \in blocks

\* Nombres únicos por directorio.
NameUniqueness ==
    \A d \in dirs :
        \A e1, e2 \in dirEntries[d] :
            (e1 /= e2) => (e1[1] /= e2[1])
\end{lstlisting}

TLC verifica las tres invariantes en cada estado alcanzable; si
una operación las rompiera, el contraejemplo lo señala. La
disciplina DPO hace que la mayoría de invariantes se mantengan
\emph{por construcción}: cada regla incluye su gluing condition
como precondición, así que TLC nunca llega a un estado donde la
invariante esté rota.

\section{Trampa: cycle por hash collision}

\begin{trampa}[BLAKE3 collision (no observada en práctica)]
La invariante~\ref{inv:sotfs-acyclic} asume que los hashes son
únicos. BLAKE3 con 256 bits tiene colisiones con probabilidad
$2^{-128}$ por par; la probabilidad de un ciclo accidental es
\emph{efectivamente cero}. Pero un atacante con poder de
construir \emph{collisions intencionales} (problema abierto en
BLAKE3) podría engañar al filesystem para que un nodo se
auto-referencie a través de su contenido.

Fix preventivo: chequeo runtime que un nodo recién insertado no
contiene su propio hash entre sus referencias. Costo: $O(\text{children})$
por inserción. Reductible si hace falta.

\textbf{Lección:} las propiedades \emph{cripto-derivadas} son
fuertes pero no infalibles. Un check runtime trivial (anti-self-reference)
preserva la invariante incluso bajo collision attack.
\end{trampa}

\section{Trampa: orphan block tras unlink}

\begin{trampa}[\emph{Memory leak} de bloques]
\texttt{unlink} borra la arista \texttt{named-as} pero el inode
puede seguir vivo si tenía $\geq 1$ \emph{otra} referencia
(typical hard-link). Si el código no rastrea el refcount
correctamente, un bloque queda \emph{orphan} (no alcanzable
desde ningún root) y nunca se libera.

Fix: cada operación que reduce el refcount marca el bloque como
\emph{candidato a GC}; el GC corre periódicamente y libera los
no alcanzables desde root. \texttt{sotFS} implementa el GC como
mark-and-sweep desde los roots de snapshot.

\textbf{Lección:} un Merkle DAG con sharing requiere GC, no
basta con refcount. Los refcounts pueden divergir bajo crash
recovery (cap.~\ref{cap:sotfs-tx}); el GC es la fuente única de
verdad.
\end{trampa}

\section{Origen del diseño}

Merkle DAGs vienen de Merkle 1987 \cite{Merkle1987-DAG}; aplicarlos
a un filesystem es de IPFS y de Git (los dos del mismo año, 2014).
El uso de DPO graph rewriting para modelar POSIX es contribución
original de \texttt{sotX} (los specs en \texttt{formal/sotfs\_*.tla}
fueron escritos antes que el código de
\texttt{services/sotfs/}). La inspiración en Yggdrasil
\cite{Sigurbjarnarson2016-Yggdrasil} para el modelo de \emph{crash
refinement} viene del cap.~\ref{cap:sotfs-tx}.

\section*{Resumen del capítulo}

\begin{itemize}
  \item sotFS modela el filesystem como un \emph{Merkle DAG}:
    nodos identificados por hash de su contenido.
  \item Seis tipos de nodo (\texttt{Inode}, \texttt{Dir},
    \texttt{Block}, \texttt{Snapshot}, \texttt{Capability},
    \texttt{Domain}); seis tipos de arista (\texttt{contains},
    \texttt{extent}, \texttt{named-as}, \texttt{parent},
    \texttt{captured}, \texttt{owns}).
  \item Direccionamiento por hash da deduplicación
    intrínseca, snapshots gratis, sharing entre dominios.
  \item Operaciones POSIX se modelan como reglas DPO con gluing
    conditions; la verificación TLA (\citaTLA{sotfs\_graph.tla})
    chequea las invariantes en cada estado.
  \item Inmutabilidad de nodos: la mutación lógica crea un nodo
    nuevo y actualiza el padre; las snapshots preservan la
    vista vieja sin copiar.
\end{itemize}

\section*{Ejercicios}

\begin{ejercicio}[$\star$]
Dibuje el grafo después de las siguientes operaciones partiendo
de un FS vacío: \texttt{mkdir /a}, \texttt{create /a/file1
content="hola"}, \texttt{create /a/file2 content="hola"}. ¿Cuántos
\texttt{Block} nodes hay? ¿Por qué?
\end{ejercicio}

\begin{ejercicio}[$\star\star$]
Lea la regla \texttt{LinkInode} del \texttt{.tla}. ¿Qué pasa si
\texttt{i \textbackslash notin inodes} (inode no vivo)? ¿Cómo
se compara con el comportamiento POSIX de \texttt{link} sobre
un inode borrado?
\end{ejercicio}

\begin{ejercicio}[$\star\star$]
Diseñe la regla DPO para \texttt{rmdir(d)}. ¿Cuál es la gluing
condition? ¿Qué ocurre si el directorio no está vacío? Compare
con la rule de \texttt{unlink}.
\end{ejercicio}

\begin{ejercicio}[$\star\star\star$]
Snapshots se hacen con un puntero a la root del DAG en un
instante. ¿Cuál es el costo en bytes de un snapshot de un FS de
1 TB con 1M de archivos? Compare con el costo de un snapshot
COW de un ext4 hipotético del mismo tamaño.
\end{ejercicio}

\begin{ejercicio}[$\star\star\star$]
Diseñe un mecanismo de \emph{garbage collection} para sotFS
basado en mark-and-sweep desde los snapshot roots. ¿Cómo
manejas concurrent writes que crean nuevos nodos durante el
mark? Pista: tri-color marking + write barrier sobre roots.
\end{ejercicio}

% TODO(apendice-E): respuestas a ejercicios 17.1-17.5
